
\chapter{\LaTeX 学术写作实践} \label{chap:methods}

本章通过具体示例介绍如何在学位论文中使用 \LaTeX 进行学术写作，包括图片插入、数学公式、代码展示、表格制作和交叉引用等常用功能。

\section{图片插入}

图片可以通过 \texttt{figure} 环境插入。以下示例插入学校的校徽图像：

\begin{figure}[htbp]
    \centering
    \includegraphics[width=4cm]{assets/nku_logo.pdf}
    \caption{南开大学校徽}
    \label{fig:nku-logo}
\end{figure}

在文中可以通过 \texttt{\textbackslash ref} 命令引用图片，如图 \ref{fig:nku-logo} 所示。

\section{数学公式}

\subsection{行内公式}

行内公式使用 \texttt{\$...\$} 或 \texttt{\textbackslash(...\textbackslash)} 包围。例如，圆的面积公式为 $A = \pi r^2$，其中 $r$ 表示半径。

\subsection{独立公式}

较复杂的公式通常独立成行，使用 \texttt{equation} 或 \texttt{align} 环境。下面是一个二次方程求根公式的示例：

\begin{equation}
    x = \frac{-b \pm \sqrt{b^2 - 4ac}}{2a}
    \label{eq:quadratic}
\end{equation}

如公式 \eqref{eq:quadratic} 所示。

\subsection{多行公式}

当公式过长时，使用 \texttt{align} 环境：
\begin{align}
    f(x) &= x^2 + 2x + 1 \\
    &= (x+1)^2
    \label{eq:expansion}
\end{align}

\section{定理与证明}

模板预置了常用的定理环境，可以直接用于学术写作：

\begin{theorem}
设函数 $f(x)$ 在区间 $[a,b]$ 上连续，则 $f(x)$ 在该区间上存在最大值与最小值。
\end{theorem}

\begin{proof}
该结论是微积分中的基本定理之一。由于闭区间上的连续函数必有界，且能够取得其上确界与下确界，因此最大值与最小值均存在。
\end{proof}

\section{代码展示}

以下是一个 Python 代码块示例：

\begin{lstlisting}[language=Python, caption={Python 代码示例}, label=code:python]
import numpy as np

x = np.linspace(0, 2*np.pi, 100)
y = np.sin(x)

print(f"数据点数: {len(x)}")
\end{lstlisting}

\section{表格制作}

使用 \texttt{booktabs} 宏包创建美观的表格：

\begin{table}[htbp]
    \centering
    \caption{实验结果对比表}
    \label{tab:results}
    \begin{tabular}{lcccc}
        \toprule
        方法 & 精度(\%) & 召回率(\%) & F1分数 & 时间(s) \\
        \midrule
        基准方法A & 88.2 & 86.5 & 87.3 & 5.2 \\
        基准方法B & 91.3 & 90.1 & 90.7 & 4.1 \\
        本论文方法 & \textbf{95.2} & \textbf{94.6} & \textbf{94.9} & \textbf{3.1} \\
        \bottomrule
    \end{tabular}
\end{table}

\section{交叉引用}

\LaTeX 提供强大的交叉引用功能：

\begin{itemize}
    \item 图片：如图 \ref{fig:nku-logo} 所示
    \item 表格：如表 \ref{tab:results} 所示
    \item 公式：如公式 \eqref{eq:expansion} 所示
    \item 代码：见代码块 \ref{code:python}
\end{itemize}

\section{参考文献引用}

在 \texttt{main.bib} 中定义参考文献，在文中使用 \texttt{\textbackslash cite} 引用。例如，参考文献 \cite{Goodfellow2016,LeCun2015} 展示了深度学习的相关工作。

\section{本章小结}

本章通过具体示例展示了图片、公式、定理、代码、表格与交叉引用等常见写作元素。熟悉这些基本功能后，大部分论文正文内容都可以在模板中较为顺畅地组织起来。
